\paragraph{\textit{Bevy}}
\todo[inline]{Uddyb om Bevy, ECS-paradigmet, herunder velegnethed og anvendelse, og spilmotor}

\paragraph{Rust}
\todo[inline]{Skriv om fordelene ved Rust omend det ikke er ``software''}

\paragraph{Photopea}
\todo[inline]{Skriv om Photopea og anvendelse heraf}
Photopea er en webapp, som forsøger at replikere det meget succefulde og kendte Photoshop. Vi valgte at benytte Photopea, da vi tidligere har erfaring med Photoshop, men det er desvære ikke kompatibelt med vores nuværende styresystem. Photopea er som førnævnt en webapp, hvilket også kommer med signe begrænsninger, herunder performance-begrænsninger. 

Vi har benyttet Photopea til at producere alle Sprites til spillet.


\paragraph{Doom Emacs}
\todo[inline]{Skriv om Doom Emacs, og anvendelse}

\paragraph{GitHub}
\todo[inline]{Skriv om GitHub repositories}
GitHub er en gratis webbaseret platform, som kan gemme Git-depoter. Git er et versionshåndteringsværktøj, som er meget udbredt i programmeringsverdenen. Med Git kan vi gemme forskellige versioner af vores kode, følge ændringer over tid og samarbejde om det samme projekt uden at overskrive hinandens arbejde. Dette gør det lettere at udvikle større programmer (som dette), da man altid kan gå tilbage til tidligere versioner, hvis der opstår fejl. GitHub fungerer samtidig som en online backup og giver mulighed for at dele projekter med andre udviklere samt udgive den endelige version af spillet så det let kan nedhentes.

GitHub tilbyder også andre værktøjer end Git. Vi har også benyttet GitHubs SCRUM-funktion samt udarbejdet et Roadmap.
